\subsection*{Configuración inicial}
\addcontentsline{toc}{subsection}{Configuración inicial}

\us[
	identificador=US-CI-01,
	titulo={Creación del repositorio de código},
	como={Desarrollador},
	quiero={Controlar las versiones de mi código, compartirlo con otros compañeros y mantener copias de seguridad para mejorar el acceso al mismo},
	requerimientos={\reqref{RNF-CA-01}, \reqref{RNF-DE-02}},
	puntos={1},
	criterios={
			\begin{itemize}
				\item El repositorio de código en la nube existe.
				\item Es accesible por todos los desarrolladores.
				\item Incluye instrucciones de uso.
			\end{itemize}
		},
	tareas={
			\begin{itemize}
				\item Inicialización de Git en un directorio local.
				\item Creación de un repositorio en Github.
				\item Configuración de acceso remoto al repositorio de Github desde la máquina local.
				\item Redacción de las instrucciones de uso en un README.
			\end{itemize}
		}
]

\us[
	identificador=US-CI-02,
	titulo={Creación de una SPA usando Vite y React},
	como={Desarrollador},
	quiero={Ver mi proyecto en funcionamiento para poder verificar que los cambios que realizo tienen el efecto deseado},
	requerimientos={\reqref{RNF-CO-01}, \reqref{RNF-DE-04}, \reqref{RNF-RE-01}},
	puntos={1},
	criterios={
			\begin{itemize}
				\item Existe una aplicación Vite + React.
				\item La aplicación puede ejecutarse localmente.
				\item La aplicación puede visitarse en el navegador.
				\item La aplicación responde de forma automática a los cambios llevados a cabo en el
				      código fuente.
			\end{itemize}
		},
	tareas={
			\begin{itemize}
				\item Creación de un proyecto Vite + React.
				\item Personalización de la configuración.
			\end{itemize}
		}
]

\us[
	identificador=US-CI-03,
	titulo={Controles de calidad con ESLint y Prettier},
	como={Desarrollador},
	quiero={Asegurar la calidad del código y mantener un estilo de codificación consistente},
	requerimientos={\reqref{RNF-CA-01}, \reqref{RNF-CA-02}},
	puntos={3},
	criterios={
			\begin{itemize}
				\item Existe una Github Action que ejecuta ESLint y Prettier sobre el código fuente.
				\item El repositorio sólo permite cambios en la rama main mediante Pull Request.
				\item La Github Action se ejecuta automáticamente cuando se crea un Pull Request.
				\item La Pull Request no se puede fusionar hasta que la Github Action haya finalizado
				      correctamente.
				\item Existen comandos documentados para ejecutar ESLint y Prettier sobre el código
				      fuente en local.
			\end{itemize}
		},
	tareas={
			\begin{itemize}
				\item Instalación de ESLint y Prettier en el proyecto.
				\item Configuración de ESLint y Prettier.
				\item Configuración de Github Actions para ejecutar ESLint y Prettier.
				\item Documentación de los comandos para ejecutar ESLint y Prettier en local.
				\item Configuración del repositorio de Github para que sólo se permitan cambios en la
				      rama main mediante Pull Request.
				\item Configuración del repositorio para hacer la Github Action obligatoria para
				      fusionar un Pull Request.
			\end{itemize}
		}
]

\us[
	identificador=US-CI-04,
	titulo={Corrección automática del código de front-end},
	como={Desarrollador},
	quiero={Tener correctores automáticos que modifican mi código para adaptarse a los estándares establecidos},
	requerimientos={\reqref{RNF-CA-02}},
	puntos={5},
	criterios={
			\begin{itemize}
				\item Existe un pre-commit hook que ejecuta ESLint y Prettier sobre el código fuente,
				      modificando el código fuente para que cumpla con los estándares establecidos.
				\item El pre-commit hook sólo se ejecuta sobre los archivos que han sido modificados
				      en el commit.
				\item El pre-commit hook corrige automáticamente los errores de formato y buenas
				      prácticas de código.
				\item El pre-commit hook muestra un informe del resultado de su ejecución.
				\item Existe documentación clara sobre el uso del pre-commit hook.
			\end{itemize}
		},
	tareas={
			\begin{itemize}
				\item Encontrar un gestor de pre-commit hooks para el proyecto, teniendo en cuenta
				      que se trata de un monorepo que tendrá sistemas en diferentes lenguajes de
				      programación.
				\item Instalar el gestor de pre-commit hooks elegido.
				\item Añadir los correctores de ESLint y Prettier.
				\item Asegurarse de que los correctores sólo se ejecutan sobre ficheros del
				      front-end.
				\item Asegurarse de que los correctores sólo se ejecutan sobre los ficheros que han
				      sido modificados en el commit.
				\item Asegurarse de que el pre-commit hook corrige automáticamente los errores de
				      formato y buenas prácticas de código.
				\item Asegurarse de que el pre-commit hook muestra un informe del resultado de su
				      ejecución.
				\item Documentación de los comandos para ejecutar ESLint y Prettier en local.
			\end{itemize}
		}
]

\us[
	identificador=US-CI-05,
	titulo={Control de ortografía y gramática},
	como={Desarrollador},
	quiero={Tener control de ortografía y gramática en los textos de la aplicación},
	requerimientos={\reqref{RNF-CA-04}},
	puntos={5},
	criterios={
			\begin{itemize}
				\item El pre-commit hook ejecuta un análisis de ortografía y gramática sobre los
				      textos visibles de la aplicación.
			\end{itemize}
		},
	tareas={
			\begin{itemize}
				\item Seleccionar un servicio de análisis de ortografía y gramática.
				\item Configurar el servicio de análisis de ortografía y gramática para que se
				      ejecute en el pre-commit hook.
				\item Documentar el uso del servicio de análisis de ortografía y gramática en el
				      pre-commit hook.
			\end{itemize}
		}
]

\us[
	identificador=US-CI-06,
	titulo={Ejecución local de la API},
	como={Desarrollador},
	quiero={Ver mi proyecto en funcionamiento para poder verificar que los cambios que realizo tienen el efecto deseado},
	requerimientos={\reqref{RNF-DE-04}},
	puntos={1},
	criterios={
			\begin{itemize}
				\item Existe una aplicación REST construida con FastAPI.
				\item La aplicación es accesible desde el entorno local.
				\item La aplicación expone un \textit{endpoint} \texttt{/docs} con su
				      documentación.
				\item La aplicación genera automáticamente su especificación OpenAPI.
				\item La aplicación expone un \textit{endpoint} \texttt{/recipes} que
				      devuelve datos fijados en el código.
				\item El servidor se reinicia automáticamente ante cambios en el código
				      fuente.
				\item Existe documentación sobre cómo lanzar el servidor en los ficheros de
				      documentación del repositorio.
			\end{itemize}
		},
	tareas={
			\begin{itemize}
				\item Creación de la API en el entorno local.
				\item Configuración de OpenAPI y del \textit{endpoint} \texttt{/docs}.
				\item Configuración del entorno local.
				\item Creación del \textit{endpoint} \texttt{/recipes}.
				\item Actualización de la documentación del repositorio.
			\end{itemize}
		}
]

\us[
	identificador=US-CI-07,
	titulo={Corrección automática del código de back-end},
	como={Desarrollador},
	quiero={Tener correctores automáticos que modifican mi código para adaptarse a los estándares establecidos},
	requerimientos={\reqref{RNF-CA-01}},
	puntos={1},
	criterios={
			\begin{itemize}
				\item Existe una comprobación del estilo de código mediante un \textit{linter} en el pre-commit hook.
				\item Existe una comprobación del formato del código mediante un \textit{formatter} en el pre-commit hook.
				\item Existe una verificación estática de tipos mediante un \textit{type checker} en el pre-commit hook.
				\item El pre-commit hook sólo se ejecuta sobre los archivos que han sido modificados en el commit.
				\item El pre-commit hook corrige automáticamente los errores de formato y buenas prácticas de código.
				\item El pre-commit hook muestra un informe del resultado de su ejecución.
				\item Existe documentación clara sobre el uso del pre-commit hook.
				\item El CI ejecuta el \textit{linter}, el \textit{formatter} y el \textit{type checker} sobre el código fuente.
			\end{itemize}
		},
	tareas={
			\begin{itemize}
				\item Configuración de un \textit{linter} para el back-end.
				\item Configuración de un \textit{formatter} para el back-end.
				\item Configuración de un \textit{type checker} para el back-end.
				\item Configuración del pre-commit hook para que ejecute el \textit{linter}, el \textit{formatter} y el \textit{type checker}.
				\item Asegurarse de que el pre-commit hook sólo se ejecuta sobre los archivos que han sido modificados en el commit.
				\item Configuración del CI para que ejecute el \textit{linter}, el \textit{formatter} y el \textit{type checker} sobre el código fuente.
			\end{itemize}
		}
]

\us[
	identificador=US-CI-08,
	titulo={Conexión del front-end con el back-end},
	como={Usuario},
	quiero={Acceder a la información de la aplicación desde el front-end de forma rápida, segura y efectiva},
	requerimientos={\reqref{RNF-RE-01}, \reqref{RNF-RE-02}},
	puntos={2},
	criterios={
			\begin{itemize}
				\item El front-end puede acceder a la información de la aplicación a través de la API.
				\item El front-end tiene un sistema de caché para almacenar la información obtenida de la API y evitar peticiones innecesarias.
				\item El front-end tiene un sistema de control de errores para manejar los errores que puedan ocurrir al acceder a la API.
				\item El front-end gestiona los tiempos de espera y las respuestas lentas de la API de forma adecuada.
			\end{itemize}
		},
	tareas={
			\begin{itemize}
				\item Creación de un servicio cliente para la API en el front-end.
				\item Sincronización automática de tipos entre el front-end y la API.
				\item Implementación de un sistema de caché en el front-end.
				\item Implementación de un sistema de control de errores en el front-end.
				\item Implementación de un sistema de gestión de tiempos de espera y respuestas lentas de la API en el front-end.
			\end{itemize}
		}
]
